\subsection{Sprint 3}

Após os problemas da entrega anterior, buscamos modificar a organização da
Sprint 3 para reduzir o acúmulo de trabalho no final do período. Em vez de
liberar somente as histórias previstas para a sprint corrente, disponibilizamos
para o time todas as histórias de usuário planejadas para as Sprints 3 e 4. A
intenção era permitir que os integrantes avançassem para outras atividades assim
que concluíssem suas tarefas, além de tornar visíveis com antecedência as
dependências e o trabalho que ainda restava até o encerramento do projeto.

Nós, alunos AGES IV, também procuramos acompanhar e cobrar com maior frequência
o andamento das tarefas, as entregas e a correção dos \textit{pull requests}.
Esse acompanhamento apresentou uma melhora em relação à Sprint 2: os PRs passaram
a receber atenção mais cedo e o fluxo de integração deixou de ficar tão parado
ao longo de toda a sprint. Para acelerar o desenvolvimento das funcionalidades
planejadas, assumimos parte dos débitos técnicos que haviam permanecido da
entrega anterior, deixando os demais integrantes mais disponíveis para avançar
nas novas histórias de usuário.

Minha atuação concentrou-se principalmente na organização do trabalho e na
preparação da entrega. Retomei a tentativa de reorganizar o ClickUp, buscando
criar um novo quadro que separasse as histórias de usuário das tarefas
individuais. O objetivo era corrigir a estrutura discutida na retrospectiva da
Sprint 2, na qual as atividades de desenvolvimento ficavam escondidas como
subtarefas. Apesar do tempo dedicado a essa reorganização, não consegui chegar a
uma configuração que resolvesse definitivamente o problema e fosse incorporada
ao fluxo de todo o time.

Próximo ao encerramento da sprint, participei de uma \textit{war room} dedicada à
entrega. Durante esse período, revisei e corrigi PRs produzidos ao longo da
semana, acompanhei as tarefas ainda abertas e ajudei a concluir o que poderia ser
incluído na apresentação. No dia da entrega, também averiguei as histórias que
não haviam sido finalizadas e procurei encaminhar as pendências restantes. Esse
trabalho ainda exigiu uma mobilização intensa, mas partiu de uma situação melhor
do que a encontrada ao final da Sprint 2, especialmente quanto ao ritmo de
revisão dos PRs.

Depois da entrega, participei do levantamento das funcionalidades que ainda não
estavam representadas no backlog. Comparamos as histórias existentes com o
protótipo original e identificamos fluxos que ainda precisavam ser registrados.
A partir desse mapeamento, criei novas histórias de usuário para completar o
backlog e tornar mais claro o trabalho necessário para a última sprint. Esse
levantamento foi importante porque a abertura antecipada das histórias das
Sprints 3 e 4 somente seria útil se o backlog representasse de forma suficiente
o escopo que ainda pretendíamos desenvolver.

\medskip

\fbox{%
    \parbox{0.95\linewidth}{%
        \setlength{\parindent}{15pt}
        \itshape A Sprint 3 apresentou uma melhora real no acompanhamento do
        trabalho e na correção dos PRs. A decisão de liberar antecipadamente as
        histórias das duas últimas sprints ajudou a reduzir a espera por novas
        tarefas e deu ao time uma visão mais ampla do que ainda precisava ser
        entregue. Assumir os débitos técnicos da Sprint 2 também foi uma forma de
        nós, AGES IV, respondermos pelos problemas acumulados e evitarmos que eles
        interrompessem novamente o desenvolvimento dos demais integrantes.

        Ainda assim, a necessidade de organizar outra \textit{war room} na
        véspera mostrou que nossa gestão continuava dependente de uma mobilização
        concentrada perto da entrega. Houve melhora, mas ainda não alcançamos um
        fluxo em que tarefas, revisões e integração fossem concluídas de maneira
        regular ao longo da sprint. Na minha atuação, consegui participar de
        forma mais preventiva da revisão dos PRs e do levantamento do backlog,
        embora a tentativa de reorganizar o ClickUp não tenha produzido o
        resultado que eu esperava.

        Também percebi que liberar mais histórias não resolve, por si só, a falta
        de clareza do backlog. O mapeamento feito após a entrega revelou que
        funcionalidades presentes no protótipo ainda não haviam sido convertidas
        em histórias de usuário. Para a gestão funcionar, não bastava aumentar a
        quantidade de tarefas disponíveis. Ainda era necessário conferir se elas
        representavam o produto, se suas dependências estavam claras e se o time
        conseguia acompanhar seu andamento.
    }%
}
